%%%%%%%%%%%%%%%%%%%%%%%%%%%%%%%%%%%%%%%%%%%%%%%%%%%%%%%%%%%%%%%%%%%%%%%%%%%%%%%
% APÉNDICE B - REGISTROS DE TIEMPOS
%%%%%%%%%%%%%%%%%%%%%%%%%%%%%%%%%%%%%%%%%%%%%%%%%%%%%%%%%%%%%%%%%%%%%%%%%%%%%%%

\chapter{Registros de tiempos}
\label{ap:tablas_dosificacion}

\normalsize

Todos los tiempos se expresan en segundos.

\section{Tiempos de la harina}

\begin{table}[!htbp]
    \centering
    \caption{Tiempos desagregados del procedimiento manual de dosificación de harina}
    \label{tab:tiempos_desagregados}

    \begingroup
    \fontsize{10}{11}\selectfont
    \setstretch{1}
    \renewcommand{\arraystretch}{1.08}

    \begin{tabular*}{\textwidth}{
        @{\extracolsep{\fill}}
        l
        c c c c c c c c c c c c
        @{}
    }

        \toprule

        \textbf{Lote} &
        \textbf{P} &
        \textbf{I} &
        \textbf{C1} &
        \textbf{R1} &
        \textbf{C2} &
        \textbf{R2} &
        \textbf{C3} &
        \textbf{R3} &
        \textbf{C4} &
        \textbf{F} &
        \textbf{D} &
        \textbf{Total} \\

        \midrule

        \multicolumn{13}{c}{
            \textbf{Jornada ordinaria --- 02-09-2026}
        } \\[0.25em]

        D1-C1 & 7,2 & 5,0 & 7,4 & 4,9 & 7,5 & 5,1 & 7,1 & --  & --  & 4,7 & 5,0 & 53,9 \\
        D1-M1 & 6,8 & 4,8 & 7,1 & 5,0 & 7,3 & 5,2 & 7,4 & 5,1 & 6,9 & 4,3 & 5,0 & 64,9 \\
        D1-F1 & 7,4 & 5,1 & 7,5 & 5,2 & 7,2 & 4,9 & 7,4 & 5,0 & 7,0 & 5,1 & 4,8 & 66,6 \\
        D1-C2 & 6,9 & 4,9 & 7,2 & 4,8 & 7,4 & 5,0 & 7,0 & --  & --  & 4,9 & 5,1 & 53,2 \\
        D1-M2 & 7,1 & 5,2 & 7,4 & 5,1 & 7,5 & 4,9 & 7,2 & 5,0 & 7,1 & 4,5 & 4,9 & 65,9 \\[0.25em]

        \multicolumn{12}{r}{\textbf{Total jornada ordinaria}} &
        \textbf{304,5} \\

        \addlinespace[0.45em]
        \midrule
        \addlinespace[0.2em]

        \multicolumn{13}{c}{
            \textbf{Jornada de sábado --- 05-09-2026}
        } \\[0.25em]

        S-C1 & 7,3 & 4,9 & 7,5 & 5,0 & 7,4 & 4,9 & 7,3 & -- & -- & 4,9 & 5,0 & 54,2 \\
        S-E1 & 6,9 & 5,0 & 7,4 & 4,9 & 7,3 & 5,0 & 7,2 & -- & -- & 4,6 & 4,9 & 53,2 \\[0.25em]

        \multicolumn{12}{r}{\textbf{Total sábado}} &
        \textbf{107,4} \\

        \addlinespace[0.2em]

        \multicolumn{12}{r}{\textbf{Total general}} &
        \textbf{411,9} \\

        \bottomrule

    \end{tabular*}

    \endgroup

    \vspace{0.3em}

    {\fontsize{10}{11}\selectfont Fuente: Elaboración propia.}

\end{table}

La Tabla~\ref{tab:tiempos_desagregados} utiliza las siguientes
abreviaturas: P, preparar el recipiente receptor, tarar y confirmar el
cero; I, desplazamiento inicial hasta la harina; C1 a C4, recoger,
trasladar y descargar cada carga principal; R1 a R3, retornar para buscar
la siguiente carga; F, realizar el ajuste y comprobar la lectura final; D,
devolver y dejar el recipiente de transporte. El símbolo ``--'' indica una
actividad no realizada.

%%%%%%%%%%%%%%%%%%%%%%%%%%%%%%%%%%%%%%%%%%%%%%%%%%%%%%%%%%%%%%%%%%%%%%%%%%%%%%%
% TIEMPOS DE AZÚCAR, SAL Y LEVADURA
%%%%%%%%%%%%%%%%%%%%%%%%%%%%%%%%%%%%%%%%%%%%%%%%%%%%%%%%%%%%%%%%%%%%%%%%%%%%%%%

\section{Tiempos de azúcar, sal y levadura}

La Tabla~\ref{tab:tiempos_secos} presenta los tiempos de las acciones del
operario durante la incorporación del azúcar, la sal y la levadura en
polvo. Entre un ingrediente y otro no se registraron pausas ni
desplazamientos, porque el operario conservó el envase en la mano.

\begin{table}[H]
    \centering
    \caption{Tiempos desagregados de la dosificación manual de azúcar, sal
    y levadura}
    \label{tab:tiempos_secos}

    \begingroup
    \fontsize{10}{11}\selectfont
    \setstretch{1}
    \renewcommand{\arraystretch}{1.08}

    \begin{tabular*}{\textwidth}{
        @{\extracolsep{\fill}}
        l
        cccc cc ccc c
        @{}
    }

        \toprule

        \multirow{2}{*}{\textbf{Lote}} &
        \multicolumn{4}{c}{\textbf{Azúcar}} &
        \multicolumn{2}{c}{\textbf{Sal}} &
        \multicolumn{3}{c}{\textbf{Levadura}} &
        \multirow{2}{*}{\textbf{Total}} \\
        \cmidrule(lr){2-5}
        \cmidrule(lr){6-7}
        \cmidrule(lr){8-10}

        &
        \textbf{I} & \textbf{L} & \textbf{R} & \textbf{V} &
        \textbf{L} & \textbf{V} &
        \textbf{L} & \textbf{V} & \textbf{E} & \\

        \midrule

        \multicolumn{11}{c}{
            \textbf{Jornada ordinaria --- 02-09-2026}
        } \\[0.25em]

        D1-C1 & 5,1 & 2,1 & 3,2 & 2,2 & 2,3 & 1,6 & 3,1 & 1,4 & 2,0 & 23,0 \\
        D1-M1 & 4,9 & 1,7 & 3,0 & 1,5 & 2,8 & 1,7 & 3,0 & 1,5 & 2,1 & 22,2 \\
        D1-F1 & 5,0 & 1,6 & 3,4 & 1,6 & 2,7 & 1,6 & 3,2 & 1,3 & 1,9 & 22,3 \\
        D1-C2 & 5,2 & 2,2 & 3,3 & 2,3 & 2,2 & 1,5 & 2,9 & 1,4 & 2,0 & 23,0 \\
        D1-M2 & 4,8 & 1,8 & 2,9 & 1,4 & 2,9 & 1,7 & 3,1 & 1,5 & 2,2 & 22,3 \\[0.25em]

        \multicolumn{10}{r}{\textbf{Total jornada ordinaria}} &
        \textbf{112,8} \\

        \addlinespace[0.45em]
        \midrule
        \addlinespace[0.2em]

        \multicolumn{11}{c}{
            \textbf{Jornada de sábado --- 05-09-2026}
        } \\[0.25em]

        S-C1 & 5,0 & 2,0 & 3,5 & 2,1 & 2,3 & 1,6 & 3,0 & 1,4 & 2,0 & 22,9 \\
        S-E1 & 4,9 & 1,9 & 3,1 & 1,9 & 2,4 & 1,5 & 2,9 & 1,3 & 1,9 & 21,8 \\[0.25em]

        \multicolumn{10}{r}{\textbf{Total sábado}} &
        \textbf{44,7} \\

        \addlinespace[0.2em]

        \multicolumn{10}{r}{\textbf{Total general}} &
        \textbf{157,5} \\

        \bottomrule

    \end{tabular*}

    \endgroup

    \vspace{0.3em}

    {\fontsize{10}{11}\selectfont Fuente: Elaboración propia.}

\end{table}

La Tabla~\ref{tab:tiempos_secos} utiliza las siguientes abreviaturas: I,
desplazamiento con el envase vacío desde la mesa hasta el saco de azúcar,
en la zona de almacenamiento; L, llenar el envase, desde el saco en el
azúcar, desde la bolsa de la mesa en la sal, e incluida la apertura del
paquete en la levadura; R, retorno directo al recipiente receptor con el
envase de azúcar; V, verter el contenido del envase sobre la harina; E,
cerrar el paquete de levadura y dejar el envase sobre la mesa.

El retorno R resultó más corto que el desplazamiento I: el operario se
dirigió directamente al recipiente receptor, mientras que I incluye tomar
el envase vacío de la mesa e iniciar el recorrido. Los desplazamientos I y
R sumaron $57{,}3\ \mathrm{s}$, el $36{,}4\%$ del tiempo de
incorporación del azúcar, la sal y la levadura.

%%%%%%%%%%%%%%%%%%%%%%%%%%%%%%%%%%%%%%%%%%%%%%%%%%%%%%%%%%%%%%%%%%%%%%%%%%%%%%%
% PROYECCIONES POR PERÍODO
%%%%%%%%%%%%%%%%%%%%%%%%%%%%%%%%%%%%%%%%%%%%%%%%%%%%%%%%%%%%%%%%%%%%%%%%%%%%%%%

\section{Proyecciones por período}

La Tabla~\ref{tab:proyecciones_periodo} presenta el tiempo y la diferencia
proyectados para distintos períodos con el mismo procedimiento de cálculo
de la Subsección~\ref{subsec:antecedentes_problema}: el número de jornadas
ordinarias y de sábados del período multiplica los valores registrados en
cada tipo de jornada. Para el azúcar, la sal y la levadura se utilizaron
las diferencias netas por jornada:
$+12$ y $+18\ \mathrm{g}$ de azúcar, $+10$ y $+4\ \mathrm{g}$ de sal, y
$+2$ y $+8\ \mathrm{g}$ de levadura, en la jornada ordinaria y en el
sábado, respectivamente. Todos los valores son cálculos derivados de las
jornadas observadas.

\begin{table}[H]
    \centering
    \caption{Tiempo y diferencia proyectados por período}
    \label{tab:proyecciones_periodo}

    \begingroup
    \footnotesize
    \setstretch{1}
    \renewcommand{\arraystretch}{1.25}
    \setlength{\tabcolsep}{3pt}

    \begin{tabular*}{\textwidth}{
        @{\extracolsep{\fill}}
        l ccc cccc
        @{}
    }
        \toprule

        \multirow{2}{*}{\textbf{\shortstack[l]{Período\\supuesto}}} &
        \multicolumn{3}{c}{\textbf{Tiempo de dosificación}} &
        \multicolumn{4}{c}{\textbf{Diferencia neta}} \\
        \cmidrule(lr){2-4}
        \cmidrule(lr){5-8}

        &
        \textbf{Harina} &
        \textbf{\shortstack{Azúcar, sal\\y levadura}} &
        \textbf{Total} &
        \textbf{\shortstack{Harina\\(kg)}} &
        \textbf{\shortstack{Azúcar\\(g)}} &
        \textbf{\shortstack{Sal\\(g)}} &
        \textbf{\shortstack{Levadura\\(g)}} \\

        \midrule

        5 ordinarias &
        25 min 22,5 s & 9 min 24,0 s & 34 min 46,5 s &
        $+0{,}390$ & $+60$ & $+50$ & $+10$ \\

        \makecell[l]{5 ordinarias\\+ 1 sábado} &
        27 min 9,9 s & 10 min 8,7 s & 37 min 18,6 s &
        $+0{,}455$ & $+78$ & $+54$ & $+18$ \\

        20 ordinarias &
        1 h 41 min 30,0 s & 37 min 36,0 s & 2 h 19 min 6,0 s &
        $+1{,}560$ & $+240$ & $+200$ & $+40$ \\

        \makecell[l]{20 ordinarias\\+ 4 sábados} &
        1 h 48 min 39,6 s & 40 min 34,8 s & 2 h 29 min 14,4 s &
        $+1{,}820$ & $+312$ & $+216$ & $+72$ \\

        \makecell[l]{22 ordinarias\\+ 4 sábados} &
        1 h 58 min 48,6 s & 44 min 20,4 s & 2 h 43 min 9,0 s &
        $+1{,}976$ & $+336$ & $+236$ & $+76$ \\

        \bottomrule
    \end{tabular*}

    \endgroup

    \vspace{0.3em}
    {\footnotesize Fuente: Elaboración propia.}
\end{table}

La diferencia neta del azúcar, la sal y la levadura corresponde a un
balance en el que los excesos y los déficits de distintos lotes se
compensan. No representa la diferencia que recibe cada preparación; para
ese fin se utilizan las diferencias por lote de la
Tabla~\ref{tab:lecturas_secos}.